\chapter*{Resumen \normalfont{{\serifthick Laburpena} Abstract}}
\addcontentsline{toc}{chapter}{Resumen \normalfont{{\serifthick Laburpena} Abstract}}

Este trabajo aborda la conversión de las actas del Pleno del Ayuntamiento de
Bilbao (2007--2026), publicadas en PDF, en dos sistemas de recuperación de
información consultables en lenguaje natural: un RAG vectorial y un sistema
GraphRAG basado en un grafo de conocimiento RDF/OWL construido y enriquecido
mediante modelos de lenguaje. El grafo modela proposiciones, grupos políticos,
concejales y temas, e incorpora razonamiento automático para responder preguntas
de agregación sobre casi 3.900 proposiciones. Ambos sistemas citan la página
exacta del documento de origen de cada dato mostrado, y se exponen mediante una
interfaz web conversacional. El trabajo incluye una evaluación comparativa
sistemática de ambas arquitecturas sobre un conjunto de preguntas de referencia
con respuesta verificada a mano.

\textbf{Palabras Clave:} GraphRAG, RAG, LLM, grafo de conocimiento, SPARQL.

%%%%%%%%%%%%%%%%%%%%%%%%%%%%%%%%%%%

\textcolor{eus-color}{{\serif [PENDIENTE: traducción al euskera del resumen anterior. No se ha generado automáticamente porque una traducción de baja calidad en un documento académico oficial es peor que dejarlo pendiente -- revisar con un hablante nativo o el servicio de euskera de la universidad.]}}

\textcolor{eus-color}{{\serif \textbf{Gako-hitzak:}} GraphRAG, RAG, LLM, ezagutza-grafoa, SPARQL.}

%%%%%%%%%%%%%%%%%%%%%%%%%%%%%%%%%%%

\textcolor{eng-color}{This work addresses the conversion of the minutes of the Plenary Council of the Bilbao City Council (2007--2026), published as PDF files, into two natural-language-queryable information retrieval systems: a vector-based RAG and a GraphRAG system built on an RDF/OWL knowledge graph constructed and enriched using large language models. The graph models proposals, political groups, councillors and topics, and incorporates automated reasoning to answer aggregation questions over nearly 3,900 proposals. Both systems cite the exact source page for every piece of data shown, and are exposed through a conversational web interface. The work includes a systematic comparative evaluation of both architectures over a hand-verified reference question set.}

\textcolor{eng-color}{\textbf{Keywords:} GraphRAG, RAG, LLM, knowledge graph, SPARQL.}
